\chapter{The Edge-Cloud Continuum and Its Architectural Paradigms}
\label{chap:edge-cloud-continuum}

Every contribution of this thesis targets the same substrate:
the computing infrastructure between cloud data centres and the devices embedded in the physical world.
%
This chapter describes that substrate,
names the environments the thesis must serve on it -- \ac{IoT} ecosystems and robot swarms --
and shows why the architectural paradigms currently used to structure distributed software cannot deploy applications at the granularity,
or with the flexibility,
those environments demand.
%
Pulverisation closes the resulting gap (\cref{chap:pulverization}).
%
The same substrate raises two further questions:
how collective systems can be programmed (\cref{chap:programming-models}),
and how their placement can be revised while they run (\cref{chap:dynamic-reconfiguration}).

\section{The Edge-Cloud Continuum}
\label{sec:the-continuum}

Cloud computing's centralised model -- offloading computation to large, remote data centres~\cite{DBLP:journals/cacm/ArmbrustFGJKKLPRSZ10} -- runs into structural limits when applications require low latency,
strict data locality,
bandwidth efficiency,
or operation under intermittent connectivity~\cite{DBLP:journals/iotj/ShiCZLX16,DBLP:journals/computer/Satyanarayanan17}.
%
Early work on mobile cloud computing already framed the tension between the resource-richness of the cloud and the constraints of end devices in terms of \emph{application partitioning}:
deciding which parts of a computation should run locally and which should be offloaded~\cite{DBLP:journals/jnca/LiuASGBQ15}.
%
The \emph{edge-cloud continuum} generalises this tension: rather than a binary split between ``device'' and ``cloud'',
it denotes a spectrum of computing and networking resources distributed between the two extremes,
over which computations can be placed,
moved,
and re-placed as conditions change,
as depicted in \cref{fig:edge-cloud-continuum}.

\begin{figure}
    \centering
    \def\svgwidth{6cm}
    \input{figures/edge-cloud-continuum.pdf_tex}
    \caption{The edge-cloud continuum is a spectrum of computing and networking resources distributed between the extremes of cloud and edge, over which computations can be placed, moved, and re-placed as conditions change.}
    \label{fig:edge-cloud-continuum}
\end{figure}

\subsection{Evolution and Topology}
\label{subsec:continuum-evolution}

The continuum emerged in stages,
as computing infrastructure moved away from the cloud.
%
Cloudlets first proposed placing a single hop of resource-rich,
virtualised infrastructure between mobile devices and the cloud,
trading the latency of a wide-area link for that of a local one~\cite{DBLP:journals/pervasive/SatyanarayananBCD09}.
%
Fog computing generalised this idea into a hierarchical layer of gateways and micro data centres between data sources and the cloud,
driven by the amount and spread of \ac{IoT} traffic~\cite{DBLP:conf/sigcomm/BonomiMZA12}.
%
Edge computing pushed the same idea further,
placing computation directly at or near the devices that generate the data rather than at an intermediate layer~\cite{DBLP:journals/iotj/ShiCZLX16,DBLP:journals/computer/Satyanarayanan17}.
%
These proposals were developed largely independently,
under many related names -- cloudlet, fog, mist, cloud-to-thing continuum -- and literature reviews have only recently tried to bring them together under one definition~\cite{DBLP:journals/access/MoreschiniPLNHT22}.

Automatic deployment of distributed software has long been recognised as a distinct engineering concern,
with dedicated models,
tools,
and open problems~\cite{DBLP:journals/jss/ArcangeliBL15}.
%
More recent work treats deployment for the continuum not as a one-off placement decision,
but as a process that can be reconfigured at any time.
%
This matches an older argument in software architecture research: a system's components and connectors should be built to change while running~\cite{DBLP:conf/icse/KramerM07}.
%
Pulverisation, for instance, breaks applications into fine-grained units whose deployment can be reconfigured at run time,
for any kind of application~\cite{DBLP:journals/fi/CasadeiPPVW20,DBLP:journals/fgcs/FarabegoliPCV24}.
%
Within the narrower field of collective and self-organising systems, the same idea appears at the level of macroprogramming:
declarative, high-level specifications state where collective behaviour should run across cloud and edge~\cite{DBLP:conf/acsos/FarabegoliVC24},
and simulation methods estimate the deployment performance of collective services once they are spread across such a topology~\cite{DBLP:journals/iotj/CasadeiFPPSV22}.
%
The resulting topology is best seen not as the fixed stack of layers suggested by the cloud/fog/edge terms,
but as a loosely structured, dynamically changing graph of heterogeneous nodes,
in which the same device can act as an edge node for one application and as a relay for another.

\subsection{Characteristics and Infrastructure Heterogeneity}
\label{subsec:continuum-characteristics}

Nodes on the continuum differ radically in computational capacity,
memory,
energy budget,
and software stack.
%
At one extreme,
microcontrollers with a few kilobytes of memory run a single sensing task under a real-time operating system,
or under no operating system at all;
at the other,
cloud servers run general-purpose operating systems and draw on effectively elastic resources.
%
This spread in resources has architectural consequences:
different classes of devices are programmed through different software component models,
historically classified along dimensions such as composition mechanisms,
deployment units,
and binding times~\cite{DBLP:journals/tse/CrnkovicSVC11}.
%
A component model built for a resource-rich server,
with dynamic loading and network-transparent binding,
does not transfer to a microcontroller that has neither.

The links connecting these nodes are just as heterogeneous:
short-range radio,
cellular,
and wired connections coexist on the same continuum,
each with its own bandwidth,
latency,
and reliability profile.
%
This is a structural property of the infrastructure,
distinct from the temporal changes in connectivity discussed in \cref{subsec:network-volatility} --
even a stationary deployment over a stable topology must reconcile a sensor on a low-power radio link with a gateway on fibre.
%
For the programmer, this asymmetry appears as a partitioning choice:
the taxonomy cited in \cref{sec:the-continuum}~\cite{DBLP:journals/jnca/LiuASGBQ15} classifies offloading strategies by what is partitioned,
when,
and under which constraints.

Energy is a further axis of heterogeneity that does not track computational capacity.
%
A gateway can be mains-powered and computationally modest,
while a mobile robot or drone carries a powerful system-on-chip on a battery budget measured in minutes,
and a sensor node may run indefinitely on a harvested,
intermittent supply.
%
Placement decisions must therefore account for how long and how often a node can sustain a computation,
not only what it can compute.
%
Energy is treated as a first-class deployment concern in \cref{subsec:energy-aware-orchestration}.

\subsection{Network Volatility and Partitioning}
\label{subsec:network-volatility}

The heterogeneity discussed in \cref{subsec:continuum-characteristics} is a static property of the infrastructure.
%
The continuum also changes over time.
%
Devices move,
radio links drop,
and nodes disappear and reappear as they run out of range or out of power --
towards the edge this becomes a defining characteristic of the network.
%
This is not a new problem in networking:
delay/disruption-tolerant networking gave up on the assumption of an end-to-end path decades ago,
buffering messages at intermediate nodes and forwarding them opportunistically as contact becomes available~\cite{DBLP:conf/sigcomm/Fall03}.
%
Opportunistic networking took the same idea into mobile ad hoc settings,
where a route is not computed in advance but assembled hop by hop as devices happen to come within range of one another~\cite{DBLP:journals/cm/PelusiPC06}.

Drone swarms are an extreme case of this volatility.
%
A \ac{FANET} is more volatile than a ground-based ad hoc network because its nodes move in three dimensions and at speed,
so links form and break faster than most routing protocols were designed for.
%
\acp{FANET} are consequently studied as a networking class of their own,
with dedicated mobility models and protocols~\cite{DBLP:journals/adhoc/BekmezciST13}.
%
Surveys of \ac{UAV} communication networks routinely list delay-tolerant,
store-carry-forward schemes as one of the standard design options,
alongside clustering and cooperative relaying~\cite{DBLP:journals/comsur/GuptaJV16}.
%
Ground robot swarms and dense \ac{IoT} deployments (\cref{subsec:swarm-robotics} and \cref{subsec:iot-systems}) face a milder version of the same problem,
since a short-range ad hoc network of many autonomous devices keeps the topology dynamic on its own.

Coordination mechanisms for the continuum need to consider such dynamic topologies as a first-class concern.
%
Time-fluid, field-based coordination, for instance, lets the collective program its own scheduling and adapt it to conditions,
instead of assuming synchrony by default~\cite{DBLP:journals/lmcs/PianiniCVMZ21}.
%
Decentralised proximity-aware clustering applies the same idea to learning:
it organises a federated learning process around whichever neighbours are reachable at a given moment,
rather than around a graph fixed in advance~\cite{DBLP:journals/iot/DominiFAVE26}.

\section{Target Environments: IoT Ecosystems and Robot Swarms}
\label{sec:target-environments}

\emph{Cyber-physical systems} (CPS)\acused{CPS} integrate computation tightly with physical processes:
they sense those processes and act back on them in feedback loops,
so their correctness depends on the physical dynamics as much as on the computation~\cite{DBLP:conf/isorc/Lee08a}.
%
\acp{CPS} and robot swarms make the heaviest demands on the continuum described in \cref{sec:the-continuum}:
they combine physical embodiment,
mobility,
and energy constraints with a need for both local autonomy and collective coordination,
so a deployment has to cope with the infrastructure's heterogeneity and its volatility together.

\subsection{Pervasive IoT Systems}
\label{subsec:iot-systems}

\ac{IoT} ecosystems~\cite{DBLP:journals/cn/AtzoriIM10} populate the continuum with large numbers of sensing and actuating devices,
connected through the heterogeneous links discussed in \cref{subsec:continuum-characteristics}.
%
Individually, such devices run comparatively simple sensing or actuation loops.
%
Coordinating many of them toward a shared goal -- monitoring a building,
a field,
or a city district -- produces the collective behaviour an \ac{IoT} deployment is built for.
%
The team of devices sustaining such a collective behaviour is not fixed:
devices join and leave as they move,
fail,
or run out of power,
so the computation must survive changes in its own set of participants (\cref{subsec:network-volatility}).
%
Programming such a changing collective as a single unit is the subject of \cref{chap:programming-models};
this section concerns the deployment problem it creates.

Any deployment of such a service must decide,
continuously,
which devices host which part of the computation.
%
Dynamic \ac{IoT} deployment reconfiguration has been studied as a global-level self-organisation problem,
in which the ecosystem itself -- rather than an external orchestrator -- decides how to redistribute computation as devices join,
leave,
or change state~\cite{DBLP:journals/iot/FarabegoliPCV24}.
%
Because trying out reconfiguration strategies on physical \ac{IoT} deployments is costly and slow,
a complementary line of work develops simulation-based methodologies to estimate the deployment performance of smart collective services before they run on real infrastructure~\cite{DBLP:journals/iotj/CasadeiFPPSV22}.

\subsection{Swarm Robotics}
\label{subsec:swarm-robotics}

Swarm robotics studies large groups of relatively simple robots that achieve complex,
useful behaviour through local interaction and self-organisation rather than central control,
a design philosophy surveyed from both a swarm-engineering and a systems perspective~\cite{DBLP:journals/swarm/BrambillaFBD13,DBLP:journals/sensors/DiasSFVCP21}.
%
Throughout this thesis, \emph{self-organisation} denotes a process in which order at the global level of a system arises solely from local interactions among its components,
without external direction or a central controller --
the sense established for biological collectives by Camazine et al.~\cite{DBLP:books/daglib/0032897}.
%
In the engineering reading of De Wolf and Holvoet,
self-organisation is distinct from emergence -- a system can exhibit either without the other -- though the two are most useful when combined~\cite{DBLP:conf/atal/WolfH04}.
%
A robot swarm is both an instance of the continuum and one of its most volatile settings (\cref{subsec:network-volatility}):
the ad hoc network connecting the robots is a side effect of their movement,
so the topology a swarm application must cope with changes as fast as the robots do.
%
Fault-tolerant cooperation among multiple robots predates this framing.
%
ALLIANCE organised a team of behaviour-based robots so that tasks could be reassigned and individual failures tolerated without central coordination~\cite{parker1998alliance},
two decades before the same requirements were attached to the continuum.

Programming a swarm directly in terms of individual robot behaviours becomes unmanageable as the swarm grows,
which motivated dedicated languages and frameworks that let a programmer express swarm behaviour at a higher level.
%
Buzz provides swarm-oriented primitives -- neighbourhood queries,
virtual stigmergy,
robot-to-robot synchronisation -- on top of a small imperative core~\cite{buzz-robot},
while Meld takes a declarative, logic-based approach and derives single-robot behaviour by compiling rules stated over the whole ensemble~\cite{meld-robot}.
%
Voltron pushes team-level programming into resource-constrained sensor networks,
targeting drones directly~\cite{voltron-robot},
and actor-based abstractions have been proposed as a general structuring principle for swarm robotic software,
isolating per-robot state behind message passing~\cite{DBLP:conf/iros/YiDLD0WY20} --
the actor model itself is treated in more general terms in \cref{subsec:actor-model-edge}.
%
Existing robotics middleware has been extended in the same direction,
as with the swarm-specific package built on top of the \ac{ROS}~\cite{rosMicros_swarm_frameworkWiki},
and toolkits such as PySwarming lower the barrier further by packaging common swarm algorithms for direct reuse~\cite{DBLP:journals/jossw/AndradeFS23}.

This body of work has been validated largely on a shared set of experimental platforms.
%
The Kilobot is a low-cost robot designed specifically to be deployed in the large numbers that swarm algorithms assume~\cite{DBLP:conf/icra/RubensteinAN12},
and the Kilogrid testbed extends it with a virtual sensing and communication substrate that lets a swarm of Kilobots be reconfigured between experiments without touching the hardware~\cite{valentini2018kilogrid}.
%
Where physical swarms of the required size are impractical,
ARGoS provides a modular, multi-engine simulation infrastructure built specifically for large swarms~\cite{pinciroli2012argos},
and Webots offers a general-purpose mobile-robot simulator with a graphical front-end~\cite{Michel2004};
both let swarm algorithms be evaluated at scale before -- or instead of -- deploying them on real robots.
%
The thesis returns to this setting at the end (\cref{chap:demonstrator}),
running one macro-program on a physical robot team to show that a deployment admitted by the thesis' model can be built on real hardware.

\subsection{Collective Adaptive Systems}
\label{subsec:collective-adaptive-systems}

\ac{IoT} ecosystems and robot swarms are best understood as two instances of a broader class,
called \emph{Collective Adaptive Systems} (CAS)\acused{CAS}:
systems composed of large numbers of autonomous, interacting entities whose collective behaviour must remain correct and useful despite continuous change in membership,
environment,
and goals~\cite{DBLP:journals/sttt/NicolaJW20}.
%
The collective \ac{IoT} services of \cref{subsec:iot-systems} and the self-organising robot teams of \cref{subsec:swarm-robotics} therefore share one requirement:
a \ac{CAS} cannot assume a fixed set of participants or a fixed topology,
so its correctness must be argued over the collective's behaviour under change,
not over any single entity in isolation.

SCEL~\cite{DBLP:journals/taas/NicolaLPT14} responds to this requirement at the language level.
%
Its communication model is attribute-based:
an interaction reaches the components whose run-time attributes satisfy a predicate,
rather than a set of recipients fixed by name,
so ensembles form and dissolve as those attributes change.
%
Membership is thus decided by local interactions alone,
making ensemble formation self-organising in the sense defined in \cref{subsec:swarm-robotics}.
%
The language comes with a formal semantics,
so the behaviour of a system can be reasoned about as its ensembles reconfigure.

Meeting this requirement calls for architecture models in which reconfiguration is part of system behaviour and has a formal account,
as already argued for self-managed systems in \cref{subsec:continuum-evolution}~\cite{DBLP:conf/icse/KramerM07}.
%
The DReAM framework instantiates this idea for \ac{CAS} specifically,
modelling dynamic reconfiguration as part of a system's architecture,
with a formal semantics that supports reasoning about the system as it reconfigures~\cite{DBLP:journals/sttt/NicolaMS20}.
%
At the programming level, a recent survey argues that \emph{macroprogramming} -- specifying the behaviour of a collective as a whole,
rather than programming individual nodes -- is the way to engineer \ac{CAS} at scale~\cite{DBLP:journals/csur/Casadei23};
the programming models that realise it are the subject of \cref{chap:programming-models}.

\subsection{Deployment and Coordination Challenges}
\label{subsec:environment-challenges}

\ac{IoT} ecosystems and robot swarms, despite their different physical form, put pressure on the continuum in the same two ways.
%
First, both must place computation and coordination logic onto nodes that are heterogeneous -- constrained sensor nodes next to gateways and cloud servers (\cref{subsec:continuum-characteristics}) --
and mobile or intermittently connected (\cref{subsec:network-volatility}),
so no placement decision can be assumed to remain valid for long.
%
Second, the coordination logic itself,
not only its placement,
must keep working as the set of participating devices changes.

Deployment for such systems has to be revisited continuously (\cref{subsec:continuum-evolution}).
%
Fine-grained deployment models such as pulverisation act on this directly:
they make the unit of deployment small enough,
and its reconfiguration cheap enough,
to be revisited as often as the environment demands~\cite{DBLP:journals/fgcs/FarabegoliPCV24}.
%
The architectural paradigms already used to structure distributed applications do not offer that much flexibility,
as shown in \cref{sec:architectural-paradigms}.

\section{Architectural Paradigms for Pervasive Systems}
\label{sec:architectural-paradigms}

Mainstream software architecture has moved through a progression of increasingly fine-grained units of deployment:
from monolithic applications,
through microservices,
to individual functions and actors.
%
Each step of this progression has been adapted to edge and pervasive settings,
since smaller units are easier to place on constrained nodes and cheaper to move between them.
%
Microservices,
\ac{FaaS},
and the actor model have each been adapted for edge deployment,
yet all three retain assumptions from their cloud origins that conflict with the fine-grained,
dynamic deployment needs of \cref{subsec:environment-challenges}.

\subsection{Microservices and FaaS at the Edge}
\label{subsec:microservices-faas-edge}

Microservices structure an application as a set of small services,
each running in its own process,
communicating through lightweight protocols,
and deployed,
scaled,
and replaced independently of the others~\cite{DBLP:books/sp/17/DragoniGLMMMS17}.
%
Independent deployability makes the model attractive for the continuum:
a unit that can be deployed on its own can,
in principle,
be placed on an edge node rather than in the cloud.
%
Osmotic computing develops this observation into an architectural vision,
in which microservices migrate between edge and cloud infrastructure in response to resource availability and application requirements,
by analogy with osmotic pressure balancing concentrations across a membrane~\cite{DBLP:journals/cloudcomp/VillariFDRR16}.

\ac{FaaS} pushes the decomposition one step further,
structuring applications as stateless functions that are triggered by events and whose provisioning,
scaling,
and lifecycle are delegated entirely to the platform~\cite{DBLP:books/sp/17/BaldiniCCCFIMMRSS17}.
%
Fine granularity and demand-driven activation suit edge workloads,
which are often bursty and tied to local events,
and dedicated platforms have followed:
serverless platforms have been prototyped specifically for edge infrastructure~\cite{DBLP:conf/icfc/BaresiM19},
and tinyFaaS shrinks the platform itself to fit a single constrained edge node,
dropping the cluster-level machinery that cloud \ac{FaaS} platforms assume~\cite{DBLP:conf/icfc/PfandzelterB20}.

The classification of software component models cited in \cref{subsec:continuum-characteristics}~\cite{DBLP:journals/tse/CrnkovicSVC11} separates the two paradigms mainly by deployment unit and binding time:
a microservice is a long-lived component bound to its host at deployment time,
whereas a function is bound at invocation and lives only for the duration of the call.
%
Both, however, assume a managed runtime substrate -- a container orchestrator or a \ac{FaaS} platform -- on every node that may host a unit,
which already restricts which parts of the continuum they can reach.

\subsection{The Actor Model at the Edge}
\label{subsec:actor-model-edge}

The actor model structures a concurrent system as a collection of actors:
isolated units of state and behaviour that interact only by asynchronous message passing,
and can create further actors in response to the messages they receive~\cite{DBLP:conf/ijcai/HewittBS73}.
%
Agha later developed the model into its canonical formal account,
in which the actor is the universal unit of both computation and distribution~\cite{DBLP:books/daglib/0066897}.
%
Because actors share no memory,
the same abstraction covers concurrency within a node and distribution across nodes,
and an actor can in principle be placed on any node able to deliver its messages.
%
The model therefore offers the finest deployment granularity of the paradigms considered here,
at the scale of an individual stateful object rather than a service or a function.

Erlang is the model's longest-standing industrial embodiment:
an Erlang node runs thousands of lightweight, isolated processes,
organised into supervision hierarchies that restart the processes that fail~\cite{DBLP:journals/cacm/Armstrong10}.
%
Decades of production use in telecommunications establish that fault isolation along actor boundaries works at scale in unreliable distributed settings.

Modern actor runtimes have made placement the platform's responsibility.
%
Orleans, for instance, activates its actors on demand,
recovers them on failure,
and reclaims them when idle,
without the application ever managing where they run~\cite{DBLP:conf/cloud/BykovGKLPT11}.
%
This location transparency decouples application logic from placement more thoroughly than microservices or \ac{FaaS} do.
%
In pervasive settings the model has been applied directly,
as in the actor-based framework for swarm robotic software already discussed in \cref{subsec:swarm-robotics},
which isolates per-robot state behind message passing~\cite{DBLP:conf/iros/YiDLD0WY20}.

\subsection{Limitations for Fine-Grained, Dynamic Deployment}
\label{subsec:paradigm-limitations}

Despite these adaptations,
all three paradigms retain deployment granularities and lifecycles designed for comparatively stable cloud infrastructure.
%
A microservice's unit of deployment is a container image running under an orchestrator,
which presumes nodes capable of hosting a container runtime and a cluster whose membership changes slowly;
neither presumption holds towards the constrained,
volatile end of the continuum characterised in \cref{subsec:continuum-characteristics} and \cref{subsec:network-volatility}.
%
Re-placing a microservice is consequently an operations-level action -- rescheduling a container -- that the application itself can neither express nor trigger.

\ac{FaaS} narrows the deployment unit but constrains it in another way:
functions are stateless by construction,
so application state moves into backing services that are typically cloud-hosted,
reintroducing the wide-area round trips that motivated edge execution in the first place.
%
The actor model avoids this by keeping state inside the actor,
and runtimes such as Orleans show that placement can be automated;
the placement policy,
however,
belongs to the runtime,
is tuned for data-centre clusters,
and is not exposed to the application as something to be programmed or adapted.
%
None of the three paradigms makes the mapping from logical units to physical hosts a first-class artefact the application can program while it runs.

Pulverisation responds directly to this gap.
%
Casadei et al.~\cite{DBLP:journals/fi/CasadeiPPVW20} split each logical device of a system into a small set of independently deployable components -- sensing,
actuation,
behaviour,
state,
and communication --
so that each component can be hosted wherever the continuum offers the resources it needs.
%
The mapping from components to hosts can then be revised at run time,
declaratively and without touching the application logic~\cite{DBLP:journals/fgcs/FarabegoliPCV24}.
%
This thesis develops that direction in \cref{chap:pulverization},
after the review of the runtime and orchestration side of the same gap in \cref{chap:dynamic-reconfiguration}.
